\section*{Chapter \ref{ch:g5:water-shapes-the-land} --- Water Shapes the Land}

\begin{solution}{exo:g5:water-shapes-the-land:1}
(a) \omterm{def:g5:water-shapes-the-land:weathering}{Weathering}: the \omterm{def:g3:stones-pebbles-sand:rock}{rock} is broken where it stands. (b) \omterm{def:g5:water-shapes-the-land:erosion}{Erosion}: the \omterm{def:g3:stones-pebbles-sand:sand}{sand}
is worn away and carried off. (c) \omterm{def:g5:water-shapes-the-land:deposition}{Deposition}: the \omterm{def:g3:stones-pebbles-sand:sand}{mud} settles. (d)
\omterm{def:g5:water-shapes-the-land:weathering}{Weathering}: the root breaks the stone where it stands.
\end{solution}

\begin{solution}{exo:g5:water-shapes-the-land:2}
The \omterm{def:g5:water-shapes-the-land:source-mouth}{source} is the place where the \omterm{def:g1:landscapes-around-us:water}{river} begins (a spring, a melting
glacier, a high wet land). The \omterm{def:g5:water-shapes-the-land:source-mouth}{mouth} is the place where it ends, in the
\omterm{def:g2:water-on-earth:ocean}{sea}, a \omterm{def:g1:landscapes-around-us:water}{lake} or a bigger \omterm{def:g1:landscapes-around-us:water}{river}. A \omterm{def:g5:water-shapes-the-land:source-mouth}{tributary} is a smaller \omterm{def:g1:landscapes-around-us:water}{river} that flows
into a bigger one.
\end{solution}

\begin{solution}{exo:g5:water-shapes-the-land:3}
The \omterm{def:g5:water-shapes-the-land:source-mouth}{source}; the narrow V-shaped \omterm{def:g1:landscapes-around-us:landform}{valley}; the \omterm{def:g5:water-shapes-the-land:meander}{meanders}; the \omterm{def:g5:water-shapes-the-land:delta}{delta}.
\end{solution}

\begin{solution}{exo:g5:water-shapes-the-land:4}
(a) False: frost, water, roots, heat and cold slowly break even hard \omterm{def:g3:stones-pebbles-sand:rock}{rock}
(\omterm{def:g5:water-shapes-the-land:weathering}{weathering}). (b) True: the faster the water, the bigger the pieces it can
move. (c) False: the heaviest pieces settle first; the finest \omterm{def:g3:stones-pebbles-sand:sand}{mud} settles
last, where the water is almost still, in a \omterm{def:g1:landscapes-around-us:water}{lake} or in the \omterm{def:g2:water-on-earth:ocean}{sea}.
\end{solution}

\begin{solution}{exo:g5:water-shapes-the-land:5}
A is on the outside of the bend: the water there is fast and deep, and it
wears the bank away (\omterm{def:g5:water-shapes-the-land:erosion}{erosion}). B is on the inside of the bend: the water
there is slow, and it drops \omterm{def:g3:stones-pebbles-sand:sand}{sand} (\omterm{def:g5:water-shapes-the-land:deposition}{deposition}).
\end{solution}

\begin{solution}{exo:g5:water-shapes-the-land:6}
In the \omterm{def:g1:landscapes-around-us:landform}{mountains} the \omterm{def:g1:landscapes-around-us:water}{river} is steep and fast: it can roll large stones
along, and it carries the small grains away. Near the \omterm{def:g2:water-on-earth:ocean}{sea} the \omterm{def:g1:landscapes-around-us:water}{river} is slow:
it can no longer move large stones, only \omterm{def:g3:stones-pebbles-sand:sand}{sand} and \omterm{def:g3:stones-pebbles-sand:sand}{mud}, and it drops them
there. On the way, too, the stones have been broken and worn into smaller
grains.
\end{solution}

\begin{solution}{exo:g5:water-shapes-the-land:7}
(a) Upper course: $600 \div 6 = 100$ kilometres. Middle course:
$600 \div 2 = 300$ kilometres. (b) Lower course: $600 - 100 - 300 = 200$
kilometres. Since $3 \times 200 = 600$, it is $\frac{1}{3}$ of the \omterm{def:g1:landscapes-around-us:water}{river}.
(Check: $\frac{1}{6} + \frac{1}{2} + \frac{1}{3} = \frac{1}{6} +
\frac{3}{6} + \frac{2}{6} = \frac{6}{6}$.)
\end{solution}

\begin{solution}{exo:g5:water-shapes-the-land:8}
The pale patches are \omterm{def:g3:stones-pebbles-sand:sand}{sand} and \omterm{def:g3:stones-pebbles-sand:pebble}{gravel} laid down on the inside of the bends,
where the water is slow (\omterm{def:g5:water-shapes-the-land:deposition}{deposition}). The dark curved \omterm{def:g1:landscapes-around-us:water}{lakes} are old loops
of the \omterm{def:g1:landscapes-around-us:water}{river}: two loops grew so close that the \omterm{def:g1:landscapes-around-us:water}{river} broke through the neck
of land between them, took the short cut, and left the old loop behind as
a \omterm{def:g1:landscapes-around-us:water}{lake}.
\end{solution}

\begin{solution}{exo:g5:water-shapes-the-land:9}
The flat meadows beside the \omterm{def:g1:landscapes-around-us:water}{river} are its \omterm{def:g5:water-shapes-the-land:flood}{floodplain}: the land the \omterm{def:g1:landscapes-around-us:water}{river}
covers when it \omterm{def:g5:water-shapes-the-land:flood}{floods}. They are flat precisely because the \omterm{def:g1:landscapes-around-us:water}{river} built them
from its own \omterm{def:g3:stones-pebbles-sand:sand}{mud}. One day the \omterm{def:g1:landscapes-around-us:water}{river} will \omterm{def:g5:water-shapes-the-land:flood}{flood} again and fill the school
with water, and fast \omterm{def:g5:water-shapes-the-land:flood}{flood} water is dangerous. The school should be built
on higher ground, and the meadows kept as fields, meadows or a park.
\end{solution}

\begin{solution}{exo:g5:water-shapes-the-land:10}
(a) A year has 12 months: $12$ billion $\div\, 12 = 1$ billion tonnes a month
on average. (b) One billion is $1\,000\,000\,000$; divided by 10, it is
$100\,000\,000$: a hundred million lorry-loads.
\end{solution}

\begin{solution}{exo:g5:water-shapes-the-land:11}
One possible story. The grain is part of a \omterm{def:g3:stones-pebbles-sand:rock}{rock} high on a \omterm{def:g1:landscapes-around-us:landform}{mountain}. Frost
opens a crack around it, night after night, until a piece of \omterm{def:g3:stones-pebbles-sand:rock}{rock} breaks
off: that is \omterm{def:g5:water-shapes-the-land:weathering}{weathering}. The piece falls to the stream below. In the upper
course the water is fast and steep: it rolls the piece along (\omterm{def:g5:water-shapes-the-land:erosion}{erosion}), and
knocking against other stones it breaks into \omterm{def:g3:stones-pebbles-sand:pebble}{pebbles} and then into grains of
\omterm{def:g3:stones-pebbles-sand:sand}{sand}. In the middle course the \omterm{def:g1:landscapes-around-us:water}{river} slows; the grain rests for a while on
the sandy beach inside a \omterm{def:g5:water-shapes-the-land:meander}{meander} bend, until a \omterm{def:g5:water-shapes-the-land:flood}{flood} picks it up again. In
the lower course the \omterm{def:g1:landscapes-around-us:water}{river} is slow and wide; the grain travels with the
water to the \omterm{def:g5:water-shapes-the-land:source-mouth}{mouth}. There the water almost stops, and the grain settles
with the \omterm{def:g3:stones-pebbles-sand:sand}{mud} and \omterm{def:g3:stones-pebbles-sand:sand}{sand} of the \omterm{def:g5:water-shapes-the-land:delta}{delta}: that is \omterm{def:g5:water-shapes-the-land:deposition}{deposition}.
\end{solution}
